% !TEX root = ../main.tex
% =====================================================================
% 1. INTRODUCTION -- page budget 1.25
%
% Written LAST, so it promises exactly what Sections 5-8 deliver and
% nothing more. Order: the setting -> the question nobody asked ->
% the four findings -> what we do not claim.
% =====================================================================

\section{Introduction}
\label{sec:intro}

A time series classifier that only returns a label is of limited use where someone has to act on
its output. When a model flags an anomaly in a patient recording or in a machine's vibration trace,
the operator needs to know \emph{which part of the signal} triggered the alert.
\citet{rudin2019stop} argues that in such settings a model that is interpretable \emph{by design}
should be preferred to a black box paired with a post-hoc explainer, because only the first
guarantees that the explanation and the decision are the same computation.

For time series, the leading realisation of that idea is Multiple Instance Learning (MIL). \millet{}
\citep{early2024millet} treats a series as a bag of time steps, replaces the global average pooling
that ends every standard convolutional classifier \citep{wang2017fcn, fawaz2020inceptiontime} with a
MIL \emph{pooling head} that scores each step, and computes the class prediction from those scores.
The scores are therefore the explanation and the prediction at once. Work in this line is normally
presented the same way: a backbone, a pooling head, an accuracy table, and a faithfulness score.

This paper asks a question that presentation leaves unanswered: \emph{within this family, what
actually controls accuracy and explanation quality?} The question is answerable because the family
factorises cleanly. A model is an encoder followed by a pooling head, the two communicate through
one fixed tensor interface, and either can be replaced without touching the other. We exploit that
to run a controlled sweep -- 72 encoder $\times$ pooling-head combinations, up to 129 datasets each,
every model trained under a single identical configuration, with the \millet{} baseline re-trained
inside the same harness rather than quoted from its paper. Under that design, any difference in the
results is attributable to the component we changed.

The sweep is built around \seanet{}, a compact multi-scale depthwise-separable encoder family and
seven MIL pooling heads, each head targeting a specific limitation of \millet{}'s conjunctive
pooling (\Cref{sec:design}). \seanet{} is the \emph{instrument} of this study rather than its
result: it exists so that capacity, aggregation rule and training budget can be varied
independently. Four findings follow, and they are the contribution of this paper.

\begin{itemize}
  \item \textbf{Selection inside the aggregator, not encoder capacity, drives explanation quality}
        (\Cref{sec:selection}). Averaging only the largest $\kappa$ fraction of per-step evidence
        makes the importance map exactly zero wherever the model used nothing. Because
        $\kappa = 1$ recovers the conjunctive baseline \emph{exactly}, sweeping $\kappa$ is a
        controlled experiment on selection alone: ground-truth ranking quality rises from $0.722$ to
        $0.777$ \ndcg{}, with an optimum rather than a monotone gain, and costs about one accuracy
        point. A second, independent knob -- an attention-entropy penalty -- moves both quantities
        the same way.

  \item \textbf{Faithful interpretability does not require a large backbone, and the residual gap is
        capacity} (\Cref{sec:cheap}). A $41$\,K-parameter model beats the $424$\,K-parameter
        baseline on accuracy, AOPCR and \ndcg{} at the same time on the one dataset where
        explanation correctness can be checked against ground truth. On the 85 \ucr{} datasets the
        small models fall behind the baseline, but a wider version of the \emph{same} encoder
        overtakes it ($11.53$ against $12.60$ mean rank), which locates the effect in capacity
        rather than in architecture.

  \item \textbf{AOPCR is not comparable across training configurations} (\Cref{sec:aopcr}). Holding
        the architecture, the code, the metric implementation and the data fixed and changing only
        the training recipe moves AOPCR from $2.57$ to $13.27$. AOPCR is an area over a curve of raw
        confidence differences, so its scale follows the scale of the logits. Every cross-paper
        AOPCR comparison in this literature is therefore uninterpretable, including comparisons
        against the published numbers of the method we build on. We give a reporting protocol that
        costs nothing to adopt.

  \item \textbf{Capacity-safety is not optimisation-safety} (\Cref{sec:safety}). Five of our seven
        heads provably reduce to the conjunctive baseline for one setting of their parameters, an
        argument used routinely to claim a design ``cannot hurt''. One of them is nonetheless the
        worst head we test, by a wide margin. The mechanism is identifiable: no term in the
        objective pushes the blend parameter back towards the safe branch, so a head that
        \emph{can} represent the baseline never finds it.
\end{itemize}

\paragraph{What we do not claim.}
We do not claim state of the art on the \ucr{} archive; we are not, and \Cref{sec:cheap} reports the
margin by which. We do not claim that the encoder variants we build are individually important --
\Cref{sec:cheap} shows the best pair is not composed of the two individually best halves. And
because of \Cref{sec:aopcr}, we do not compare any AOPCR value in this paper against an AOPCR
printed in another one. Every number we report was produced under a single configuration inside one
harness, and the code that produced it is available at \anonrepo{}.
